\documentclass[11pt,a4paper]{article}
\usepackage[T2A]{fontenc}
\usepackage[utf8]{inputenc}
\usepackage[russian]{babel}
\usepackage{amsmath,amssymb}
\usepackage[left=20mm,right=20mm,top=16mm,bottom=15mm,
  includeheadfoot,headheight=15pt,headsep=7mm,footskip=8mm]{geometry}
\usepackage{xcolor,needspace,fancyhdr}
\usepackage{tcolorbox}
\tcbuselibrary{skins,breakable}
\usepackage[unicode,hidelinks]{hyperref}
\hypersetup{pdftitle={Зачёт по КР2},pdfsubject={Дифференциальные уравнения. Вопросы 1–17}}

% Палитра и накладные заголовки из основного конспекта.
\definecolor{HSEblue}{HTML}{004C97}
\definecolor{HSEgreen}{HTML}{007F5F}
\definecolor{HSEteal}{HTML}{008B8B}
\definecolor{HSEpurple}{HTML}{6E56DB}
\definecolor{HSEgray}{HTML}{666666}
\tcbset{hsebox/.style={enhanced,breakable,arc=6pt,boxrule=.8pt,
  colback=HSEgray!8,coltitle=white,fonttitle=\small\bfseries,
  left=8pt,right=8pt,top=7pt,bottom=5pt,boxsep=1.5pt,
  before skip=8pt,after skip=6pt,
  attach boxed title to top left={xshift=2mm,yshift=-2mm},
  boxed title style={frame hidden,boxrule=0pt,arc=2pt,left=4pt,right=4pt,top=1pt,bottom=1pt},
  before upper={\setlength{\parindent}{0pt}\setlength{\parskip}{3pt}}}}
\tcbset{paired/.style={sidebyside,sidebyside align=center seam,sidebyside gap=4mm,
  righthand width=.43\linewidth,lower separated=false}}
\newcommand{\declarecard}[2]{%
  \newtcolorbox{#1}[2][]{hsebox,colframe=#2!75!black,colbacktitle=#2!85!black,title={##2},##1}}
\declarecard{definition}{HSEblue}
\declarecard{result}{HSEgreen}
\declarecard{lemma}{HSEteal}
\declarecard{method}{HSEpurple}

\pagestyle{fancy}
\fancyhf{}
\fancyhead[L]{\small\color{HSEblue}\bfseries Зачёт по КР2}
\fancyhead[R]{\small\itshape Дифференциальные уравнения}
\fancyfoot[C]{\thepage}
\renewcommand{\headrulewidth}{.4pt}
\fancypagestyle{opening}{\fancyhf{}\fancyfoot[C]{\thepage}\renewcommand{\headrulewidth}{0pt}}

\setlength{\parindent}{0pt}
\setlength{\parskip}{4pt plus 1pt}
\raggedbottom
\newcommand{\question}[3]{%
  \par\addvspace{10pt}\Needspace{8\baselineskip}%
  \pdfbookmark[1]{#1. #2}{question.#1}%
  {\large\bfseries\color{HSEblue}#1.\enspace #2\par}\nobreak
  {\small\itshape\raggedright #3\par}\nobreak\smallskip}
\newcommand{\idea}[1]{%
  \par\noindent{\color{HSEteal!65!black}\bfseries Идея доказательства.}\enspace #1\par}

\begin{document}
\thispagestyle{opening}
{\Huge\bfseries\color{HSEblue}Зачёт по КР2\par}
{\large Дифференциальные уравнения\par}
\question{1}{Уравнения с разделяющимися переменными}{Сформулируйте и поясните теорему существования и единственности
для уравнений с разделяющимися переменными. Расскажите о
частных случаях}
\begin{result}{Теорема}
Пусть \(f\) непрерывна на интервале \(I\), а \(g\) непрерывна и нигде не обращается в нуль на интервале \(K\). Для любой точки \((x_0,y_0)\in I\times K\) задача \(y'=f(x)g(y)\), \(y(x_0)=y_0\), имеет единственное локальное решение. Оно определяется равенством
\[
 \int_{y_0}^{y(x)}\frac{du}{g(u)}=\int_{x_0}^{x}f(t)\,dt .
\]
\end{result}

\idea{Деление на \(g(y)\) и интегрирование дают это равенство. Его левая часть строго монотонна по верхнему пределу и имеет дифференцируемую обратную функцию.}

При \(y'=f(x)\) получаем \(y(x)=y_0+\int_{x_0}^{x}f(t)\,dt\), поэтому интегральные кривые отличаются вертикальными сдвигами. Если \(f(x)\to+\infty\) или \(f(x)\to-\infty\) у конечного края интервала, то при сходимости интеграла решение имеет конечный предел, а при расходимости возникает вертикальная асимптота.

При \(y'=g(y)\) на каждом интервале, где \(g\ne0\), имеем \(x-x_0=\int_{y_0}^{y}du/g(u)\). Кривые в такой полосе отличаются горизонтальными сдвигами. Каждый нуль \(g(c)=0\) даёт решение \(y\equiv c\) и для общего уравнения \(y'=f(x)g(y)\). В автономном случае соседняя ветвь достигает \(c\) за конечное время, если интеграл \(\int_{y_0}^{c}du/g(u)\) сходится. Склеивая её с постоянным решением, получаем нарушение единственности. При расходимости интеграла эта ветвь лишь асимптотически приближается к \(y=c\).

\question{2}{Однородная правая часть}{Сформулируйте и поясните теорему существования и единственности
для уравнений с однородной правой частью.}
\begin{result}{Теорема}
Рассмотрим задачу \(y'=f(y/x)\), \(y(x_0)=y_0\), где \(x_0\ne0\). Если \(f\) непрерывна на интервале \(K\) и \(f(u)-u\ne0\) при всех \(u\in K\), то через каждую точку с \(x\ne0,\ y/x\in K\) проходит единственное локальное решение. Если \(u_0=y_0/x_0\), его задаёт равенство
\[
 \int_{u_0}^{\,y(x)/x}\frac{du}{f(u)-u}=\ln\left|\frac{x}{x_0}\right|.
\]
\end{result}

\idea{Замена \(u=y/x\) превращает исходное уравнение в \(u'=(f(u)-u)/x\). К нему применима теорема для уравнения с разделяющимися переменными.}

Каждый корень \(f(c)=c\) даёт решение \(y=cx\) на каждой из полупрямых \(x>0\) и \(x<0\). Такие решения теряются при делении на \(f(u)-u\). Для единственности через точку такой прямой нужны дополнительные сведения о \(f\), например её локальная липшицевость вблизи \(c\).

\clearpage
\question{3}{Уравнение высшего порядка}{Сформулируйте теорему существования и единственности для
уравнения n-го порядка, разрешенного относительно старшей
производной}
\begin{result}{Теорема}
Пусть \(G\subset\mathbb R^{n+1}\) --- область, в которой непрерывны \(F(x,z_0,\ldots,z_{n-1})\) и все частные производные \(\partial F/\partial z_j\), \(0\le j<n\). Тогда для каждой начальной точки \((x_0,a_0,\ldots,a_{n-1})\in G\) существует \(h>0\), для которого задача
\[
 y^{(n)}=F(x,y,y',\ldots,y^{(n-1)}),\qquad
 y^{(j)}(x_0)=a_j,\quad j=0,\ldots,n-1,
\]
имеет единственное решение на \([x_0-h,x_0+h]\). Под \(y^{(0)}\) понимается сама функция \(y\).
\end{result}

\question{4}{Нормальная система ОДУ}{Дайте определение нормальной системы ОДУ, решения нормальной
системы ОДУ, задачи Коши для нормальной системы ОДУ,
сформулируйте теорему существования и единственности для
нормальной системы ОДУ}
\begin{definition}{Определения}
Нормальная система порядка \(n\) имеет вид \(\dot x^i=F^i(t,x^1,\ldots,x^n)\), \(i=1,\ldots,n\), или \(\dot{\mathbf x}=\mathbf F(t,\mathbf x)\). Компоненты правой части непрерывны в области \(G\subset\mathbb R^{n+1}\). Решение на промежутке \(J\) --- дифференцируемая вектор-функция \(\boldsymbol\varphi\colon J\to\mathbb R^n\), для которой при каждом \(t\in J\)
\[
 (t,\boldsymbol\varphi(t))\in G,\qquad
 \dot{\boldsymbol\varphi}(t)=\mathbf F(t,\boldsymbol\varphi(t)).
\]
Задача Коши состоит в нахождении решения с заданным значением \(\boldsymbol\varphi(t_0)=\mathbf x_0\).
\end{definition}

\begin{result}{Теорема}
Если все \(F^i\) и \(\partial F^i/\partial x^j\), \(1\le i,j\le n\), непрерывны в \(G\), то для каждой точки \((t_0,\mathbf x_0)\in G\) найдётся \(h>0\), при котором решение задачи Коши существует и единственно на \([t_0-h,t_0+h]\).
\end{result}

\question{5}{Линейная система}{Дайте определение линейно однородной и линейно неоднородной
системы. Расскажите о применении теоремы существования и
единственности к линейной системе с непрерывной матрицей и
непрерывной вектор-функцией в правой части}
\begin{definition}{Однородная и неоднородная системы}
Линейная система имеет вид \(\dot{\mathbf x}=A(t)\mathbf x+\mathbf f(t)\), где \(A(t)=(a_{ij}(t))\) --- матрица размера \(n\times n\), а \(\mathbf f(t)\) --- вектор из \(n\) компонент. При \(\mathbf f\equiv\mathbf0\) система называется линейной однородной (ЛОС), при \(\mathbf f\not\equiv\mathbf0\) --- линейной неоднородной (ЛНС).
\end{definition}

\begin{result}{Существование и единственность}
Если все компоненты \(A\) и \(\mathbf f\) непрерывны на промежутке \(I\), то для любых \(t_0\in I\) и \(\mathbf x_0\in\mathbb R^n\) задача с условием \(\mathbf x(t_0)=\mathbf x_0\) имеет единственное решение на всём \(I\). Действительно, для \(F^i(t,\mathbf x)=\sum_{j=1}^n a_{ij}(t)x^j+f^i(t)\) имеем \(\partial F^i/\partial x^j=a_{ij}(t)\). Поэтому теорема для нормальной системы даёт локальное существование и единственность. На каждом компактном отрезке внутри \(I\) коэффициенты ограничены, и оценка Гронуолла исключает уход решения на бесконечность за конечное время внутри \(I\).
\end{result}

\clearpage
\question{6}{Линейная зависимость функций}{Дайте определение линейно зависимой системы функций, линейно
независимой системы вектор-функций. Сформулируйте и поясните
утверждение о связи линейной зависимости вектор-функций и их
значений. Дайте определение вронскиана n-мерных вектор-функций.}
\begin{definition}{Линейная зависимость и независимость}
Функции \(u_1,\ldots,u_k\) линейно зависимы на \(M\), если существуют постоянные \(c_1,\ldots,c_k\), не все равные нулю, для которых \(\sum_{j=1}^k c_j u_j(t)=0\) при всех \(t\in M\). Вектор-функции \(\mathbf v_1,\ldots,\mathbf v_k\) линейно независимы на \(M\), если из тождества \(\sum_{j=1}^k c_j\mathbf v_j(t)=\mathbf0\) следует \(c_1=\cdots=c_k=0\).
\end{definition}

Из линейной зависимости вектор-функций на \(M\) следует зависимость их значений в каждой точке \(M\). Следовательно, независимости значений хотя бы в одной точке достаточно для независимости вектор-функций на \(M\).

\idea{В тождество с постоянными коэффициентами можно подставить любое \(t\in M\). При поточечной зависимости набор коэффициентов может меняться вместе с \(t\).}

Например, \(\mathbf v_1(t)=(1,0)^{\mathsf T}\) и \(\mathbf v_2(t)=(t,0)^{\mathsf T}\) коллинеарны при каждом \(t\in\mathbb R\). Но равенство \(c_1+c_2t=0\) для всех \(t\) возможно только при \(c_1=c_2=0\), так что вектор-функции независимы.

\begin{definition}{Определитель Вронского}
Для \(n\) вектор-функций \(\mathbf v_j\colon M\to\mathbb R^n\) определитель Вронского есть
\[
 W(t)=\det\bigl(\mathbf v_1(t)\ \cdots\ \mathbf v_n(t)\bigr).
\]
Столбец с номером \(j\) состоит из координат \(\mathbf v_j(t)\).
\end{definition}

\question{7}{Вронскиан решений системы}{Сформулируйте и поясните утверждения о связи равенства нулю
определителя Вронского и линейной зависимости вектор-функций.
Сформулируйте лемму о связи равенства нулю вронскиана решений
ЛОС в некоторой точке и линейной зависимости данных функций.}
Пусть \(M\) --- промежуток, \(x_1,\ldots,x_n\colon M\to\mathbb R^n\),
\(X(t)=(x_1(t)\ \cdots\ x_n(t))\) и \(W(t)=\det X(t)\).
Линейная зависимость на \(M\) означает, что найдётся постоянный вектор
\(c\ne0\), для которого \(X(t)c=0\) при всех \(t\in M\).
Тогда \(W\equiv0\). Если же \(W(t_0)\ne0\) хотя бы в одной точке,
вектор-функции линейно независимы. Для произвольных вектор-функций даже
\(W\equiv0\) ещё не влечёт зависимости на всём \(M\).

\begin{lemma}{Лемма о решениях ЛОС}
Если \(x_1,\ldots,x_n\) решают одну ЛОС \(x'=A(t)x\) с непрерывной
матрицей \(A\), то уже из \(W(t_0)=0\) в одной точке следует линейная
зависимость на всём \(M\), а значит, и \(W\equiv0\).
\end{lemma}

\idea{Зависимость функций даёт при каждом \(t\) зависимые столбцы, поэтому определитель равен нулю. Для решений ЛОС условие \(W(t_0)=0\) даёт ненулевой набор коэффициентов, с которым комбинация решений равна нулю в \(t_0\), а по единственности и на всём \(M\).}

\clearpage
\question{8}{Фундаментальная система решений}{Дайте определение фундаментальной системы решений ЛОС. Дайте
определение общего решения ЛОС. Сформулируйте и поясните
теорему об общем решении ЛОС. Дайте определение
фундаментальной матрицы ЛОС. Сформулируйте теорему о переходе
от одной фундаментальной матрицы ЛОС к другой.}
Пусть \(x'=A(t)x\) --- ЛОС порядка \(n\) с непрерывной на промежутке
\(M\) матрицей. Её \emph{фундаментальная система решений} (ФСР) ---
любые \(n\) линейно независимых на \(M\) решений. \emph{Общее решение}
--- семейство, каждый элемент которого решает систему и которое содержит
все её решения на \(M\).

\begin{result}{Теорема об общем решении}
Пространство решений ЛОС имеет размерность \(n\),
а любая ФСР \(x_1,\ldots,x_n\) образует его базис. Общее решение имеет вид
\[
 x(t)=c_1x_1(t)+\cdots+c_nx_n(t),\qquad c_1,\ldots,c_n\in\mathbb R .
\]
\end{result}

ФСР можно получить, взяв решения с начальными значениями \(x_j(t_0)=e_j\),
где \(e_j\) --- векторы стандартного базиса.

\idea{Значения решений ФСР в \(t_0\) образуют базис, поэтому начальный вектор однозначно разлагается по ним. По единственности та же комбинация решений даёт искомое решение на всём \(M\).}

\emph{Фундаментальная матрица} \(X(t)=(x_1(t)\ \cdots\ x_n(t))\) ---
матрица со столбцами из ФСР. Она удовлетворяет \(X'=AX\),
\(\det X(t)\ne0\) на \(M\), и общее решение записывается как \(x=Xc\).
\begin{result}{Переход между фундаментальными матрицами}
Если \(C\) --- постоянная невырожденная
матрица, то \(Y=XC\) также фундаментальна. Для любых двух фундаментальных
матриц одной системы существует единственная такая \(C\):
\(C=X(t_0)^{-1}Y(t_0)\).
\end{result}

\question{9}{Дифференцирование определителя}{Выпишите формулу дифференцирования определителя матрицы,
докажите её при n =2. Запишите формулу Лиувилля-Остроградского
для вронскиана решений ЛОС.}
\begin{lemma}{Дифференцирование определителя}
Пусть \(X(t)\) --- квадратная матрица с дифференцируемыми элементами,
а \(D_i(t)\) --- определитель, в котором \(i\)-я строка \(X(t)\) заменена
её производной. Тогда
\[
 \frac{d}{dt}\det X(t)=\sum_{i=1}^{n}D_i(t).
\]
\end{lemma}
\emph{Доказательство при \(n=2\).} По правилу производной произведения
\[
 \frac{d}{dt}\begin{vmatrix}a&b\\c&d\end{vmatrix}
 =(ad-bc)'=(a'd-b'c)+(ad'-bc')
 =\begin{vmatrix}a'&b'\\c&d\end{vmatrix}
  +\begin{vmatrix}a&b\\c'&d'\end{vmatrix}.
\]

\begin{result}{Формула Лиувилля--Остроградского}
Пусть столбцы \(X(t)\) ---
\(n\) решений ЛОС \(x'=A(t)x\) с непрерывной на промежутке матрицей \(A\),
\(W(t)=\det X(t)\). Тогда
\[
 W(t)=W(t_0)\exp\!\left(\int_{t_0}^{t}\operatorname{tr}A(s)\,ds\right).
\]
\end{result}

\clearpage
\question{10}{Вариация постоянных для системы}{Сформулируйте утверждение о связи общего решения ЛНС и ЛОС.
Расскажите о методе вариации постоянных решения ЛНС.}
\begin{result}{Общее решение ЛНС}
Пусть \(A(t)\) и \(f(t)\) непрерывны на промежутке \(M\), а \(X(t)\) ---
фундаментальная матрица ЛОС \(x'=A(t)x\). Общее решение ЛОС есть \(x=Xc\),
\(c\in\mathbb R^n\). Если \(x_p\) --- одно частное решение ЛНС
\(x'=A(t)x+f(t)\), то её общее решение есть \(x=x_p+Xc\):
разность любых двух решений ЛНС решает ЛОС, а прибавление решения ЛОС
к \(x_p\) снова даёт решение ЛНС.
\end{result}

\begin{method}{Вариация постоянных}
Положим \(x=X(t)c(t)\). Тогда
\(X'c+Xc'=AXc+f\). Поскольку \(X'=AX\), остаётся \(Xc'=f\), откуда
\(c'=X^{-1}f\). После интегрирования
\[
 x(t)=X(t)c_0+X(t)\int_{t_0}^{t}X(s)^{-1}f(s)\,ds,\qquad c_0\in\mathbb R^n .
\]
При \(c_0=0\) получаем частное решение с \(x(t_0)=0\).
\end{method}

\question{11}{Переход от уравнения к системе}{Расскажите о переходе от линейного уравнения n-го порядка к
линейной системе, сформулируйте утверждения о сохранении
линейной зависимости решений линейного уравнения и линейной
системы при данном переходе. Сформулируйте теорему об общем
решении линейного однородного уравнения (ЛОУ).}
Пусть коэффициенты и правая часть линейного уравнения
\[
 y^{(n)}+a_{n-1}(t)y^{(n-1)}+\cdots+a_1(t)y'+a_0(t)y=f(t)
\]
непрерывны на промежутке \(M\). Замена \(x^k=y^{(k-1)}\),
\(k=1,\ldots,n\), приводит к системе
\[
 \begin{aligned}
  (x^k)'&=x^{k+1}, &&1\le k<n,\\
  (x^n)'&=-a_0(t)x^1-\cdots-a_{n-1}(t)x^n+f(t).
 \end{aligned}
\]
Каждое решение \(y\) даёт решение системы
\(x=(y,y',\ldots,y^{(n-1)})^{\mathsf T}\). Обратно, первые \(n-1\)
равенств определяют все координаты по \(y=x^1\), а последнее даёт
исходное уравнение.

Линейная зависимость сохраняется в обе стороны. Если
\(\sum_j c_jy_j(t)\equiv0\) для постоянных \(c_j\), не всех нулевых,
то после последовательного дифференцирования до порядка \(n-1\)
получаем \(\sum_j c_jx_j(t)\equiv0\). Обратно, первая координата
последнего равенства даёт \(\sum_j c_jy_j(t)\equiv0\).

\begin{result}{Теорема об общем решении ЛОУ}
При \(f\equiv0\) решения образуют
линейное пространство размерности \(n\). Любые \(n\) независимых решений
\(y_1,\ldots,y_n\) образуют его ФСР, и каждое решение единственным образом
записывается как \(y=\sum_{j=1}^n c_jy_j(t)\), \(c_j\in\mathbb R\).
Это следует из линейного взаимно однозначного соответствия с решениями ЛОС.
\end{result}

\clearpage
\begingroup
\setlength{\abovedisplayskip}{4pt}
\setlength{\belowdisplayskip}{4pt}
\setlength{\parskip}{2pt}
\tcbset{hsebox/.append style={before skip=6pt,after skip=4pt}}
\question{12}{Вронскиан скалярных функций}{Дайте определение вронскиана системы функций и расскажите о связи
линейной зависимости функций и равенства определителя Вронского
этих функций нулю. Запишите формулу Лиувилля-Остроградского для
линейных уравнений.}
\begin{definition}[paired]{Вронскиан}
Для \(n\) функций \(y_1,\ldots,y_n\), имеющих производные до порядка
\(n-1\) на промежутке \(M\), \emph{вронскианом} называют \(W(t)=\det\!\bigl(y_j^{(i-1)}(t)\bigr)_{i,j=1}^{n}\).
\tcblower
\[ W(t)=\begin{vmatrix}
  y_1&\cdots&y_n\\ y_1'&\cdots&y_n'\\
  \vdots&&\vdots\\ y_1^{(n-1)}&\cdots&y_n^{(n-1)}
 \end{vmatrix}.
\]
\end{definition}
Линейная зависимость функций влечёт \(W\equiv0\): одно и то же
нетривиальное соотношение между столбцами получается дифференцированием
\(\sum_jc_jy_j\equiv0\). Поэтому \(W(t_0)\ne0\) хотя бы в одной точке
гарантирует независимость. Для произвольных функций равенства \(W\equiv0\) недостаточно для вывода о зависимости. Если все \(n\) функций решают одно ЛОУ
порядка \(n\) с непрерывными коэффициентами, то и обратный вывод верен:
из \(W(t_0)=0\) следует зависимость на всём \(M\). Это лемма для решений
соответствующей ЛОС, применённая к вектор-функциям
\((y_j,y_j',\ldots,y_j^{(n-1)})^{\mathsf T}\).

\begin{result}[paired]{Формула Лиувилля--Остроградского}
Для \(n\) решений уравнения
\(y^{(n)}+a_{n-1}(t)y^{(n-1)}+\cdots+a_0(t)y=0\) справедлива формула
\tcblower
\[
 W(t)=W(t_0)\exp\!\left(-\int_{t_0}^{t}a_{n-1}(s)\,ds\right).
\]
\end{result}
При переходе к системе вронскиан сохраняется, а след её матрицы равен
\(-a_{n-1}(t)\), что и объясняет знак в показателе экспоненты.

\question{13}{Построение уравнения и вариация постоянных}{Сформулируйте и поясните теорему о построении ЛОУ по известной
фундаментальной системе решений. Расскажите о методе вариации
постоянных решения линейного неоднородного уравнения(ЛНУ).}
\begin{result}[paired]{Построение уравнения по ФСР}
Пусть $y_1,\ldots,y_n\in C^n(I)$ и их вронскиан $W(x)$ нигде не равен нулю
на промежутке $I$. Тогда ЛОУ $n$-го порядка с этой ФСР задаётся уравнением справа. После деления на $W(x)$ получаем единственное уравнение с этой ФСР и единицей при $y^{(n)}$. Его коэффициенты непрерывны.
\tcblower
\[
 \begin{vmatrix}
 y_1&\cdots&y_n&y\\
 y'_1&\cdots&y'_n&y'\\
 \vdots&&\vdots&\vdots\\
 y_1^{(n)}&\cdots&y_n^{(n)}&y^{(n)}
 \end{vmatrix}=0.
\]
\end{result}

\idea{При разложении по последнему столбцу коэффициент при $y^{(n)}$ равен $W(x)$. Подстановка $y=y_j$ даёт два одинаковых столбца, а условие $W\ne0$ гарантирует независимость этих $n$ решений.}

\begin{method}{Вариация постоянных}
Для ЛНУ с единицей при $y^{(n)}$ и правой частью
$f(x)$ берём ФСР соответствующего ЛОУ и ищем $y=\sum_{j=1}^n c_j(x)y_j(x)$.
Производные $c'_j$ определяем из системы
\[
 \sum_{j=1}^n y_j^{(r)}c'_j=0\quad(0\le r\le n-2),\qquad
 \sum_{j=1}^n y_j^{(n-1)}c'_j=f(x).
\]
Определитель системы равен $W(x)\ne0$. Интегрируем $c'_j$ и подставляем $c_j$
в выражение для $y$. Произвольные постоянные интегрирования дают общее решение.
\end{method}

\endgroup
\clearpage
\question{14}{Простые корни}{Дайте определение характеристического многочлена ЛОУ.
Сформулируйте и поясните теорему об общем решении ЛОУ с
постоянными коэффициентами в случае некратных корней
характеристического уравнения.}
\begin{definition}{Характеристический многочлен}
Для ЛОУ $a_n y^{(n)}+\cdots+a_1y'+a_0y=0$ с постоянными коэффициентами
и $a_n\ne0$ \textit{характеристический многочлен} имеет вид
\[
 L(\lambda)=a_n\lambda^n+\cdots+a_1\lambda+a_0.
\]
\end{definition}
\begin{result}{Теорема о простых корнях}
Если все $n$ корней $\lambda_1,\ldots,\lambda_n$ простые, то функции
$e^{\lambda_1x},\ldots,e^{\lambda_nx}$ образуют ФСР над $\mathbb C$ и
\[
 y(x)=\sum_{j=1}^n C_j e^{\lambda_jx},\qquad C_j\in\mathbb C .
\]
При вещественных коэффициентах пару корней $\alpha\pm i\beta$ в вещественной
ФСР представляют функции $e^{\alpha x}\cos(\beta x)$ и $e^{\alpha x}\sin(\beta x)$.
\end{result}

\idea{Подстановка $e^{\lambda x}$ в левую часть даёт $L(\lambda)e^{\lambda x}$. Вронскиан при $x=0$ равен определителю Вандермонда $\prod_{i<j}(\lambda_j-\lambda_i)\ne0$, поэтому эти $n$ решений независимы.}

\question{15}{Кратные корни и формула смещения}{Сформулируйте и поясните теорему об общем решении ЛОУ с
постоянными коэффициентами в случае кратных корней
характеристического уравнения, предварительно записав формулу
смещения}
\begin{lemma}{Формула смещения}
Пусть $D=d/dx$, а $L(D)$ --- многочлен от $D$ с постоянными коэффициентами.
Для $v\in C^n(I)$, где $n=\deg L$, справедлива \textit{формула смещения}
\[
 L(D)\bigl(e^{\lambda x}v(x)\bigr)=e^{\lambda x}L(D+\lambda)v(x).
\]
\end{lemma}
\begin{result}{Теорема о кратных корнях}
Если различные корни $\lambda_1,\ldots,\lambda_m$ имеют кратности
$k_1,\ldots,k_m$, то каждому $\lambda_j$ в ФСР отвечают
$e^{\lambda_jx},xe^{\lambda_jx},\ldots,x^{k_j-1}e^{\lambda_jx}$. Общее решение
\[
 y(x)=\sum_{j=1}^m e^{\lambda_jx}P_j(x),\qquad \deg P_j\le k_j-1,
\]
где $P_j$ --- произвольные многочлены, включая нулевой.
Для вещественного уравнения пара $\alpha\pm i\beta$ кратности $k$ даёт
$x^r e^{\alpha x}\cos(\beta x)$ и $x^r e^{\alpha x}\sin(\beta x)$ при $0\le r<k$.
\end{result}

\idea{После сдвига множитель $(D-\lambda_j)^{k_j}$ превращается в $D^{k_j}$ и уничтожает $P_j$. Применяя $\prod_{i\ne j}(D-\lambda_i)^{k_i}$ к тождественному нулю, исключаем остальные группы и получаем $P_j=0$, поскольку в группе $j$ этот оператор сохраняет степень ненулевого многочлена.}

\clearpage
\question{16}{Квазимногочлен и резонанс}{Дайте определение квазимногочлена. Как связаны общее решение
ЛНУ и общее решение ЛОУ? Сформулируйте теорему о частном
решении ЛНУ с постоянными коэффициентами с правой частью в виде
квазимногочлена}
\begin{definition}{Квазимногочлен}
\textit{Квазимногочлен} --- конечная сумма функций $e^{\mu x}P(x)$, где
$\mu\in\mathbb C$ и $P$ --- многочлен. Общее решение ЛНУ $L(D)y=f$ есть сумма
одного частного решения этого уравнения и общего решения ЛОУ $L(D)y=0$.
\end{definition}

\begin{result}{Частное решение и резонанс}
Пусть коэффициенты $L(D)$ постоянны, $f(x)=e^{\mu x}P_N(x)$ и $\deg P_N=N$.
Тогда существует частное решение
\[
 y_{\text{ч}}(x)=x^s e^{\mu x} Q_N(x),\qquad \deg Q_N=N.
\]
Если $L(\mu)\ne0$, то $s=0$. При резонансе $s$ равняется кратности корня
$\mu$ характеристического многочлена $L$. Коэффициенты $Q_N$ находим
подстановкой и сравнением коэффициентов при одинаковых степенях $x$.
Для конечной суммы таких слагаемых складываем соответствующие частные решения.
\end{result}

\question{17}{Интегральная формула Коши}{Сформулируйте и поясните интегральную формулу Коши нахождения
частного решения ЛНУ.}
\begin{definition}{Ядро Коши}
Рассмотрим ЛНУ $L[y]=y^{(n)}+a_{n-1}(x)y^{(n-1)}+\cdots+a_0(x)y=f(x)$
с непрерывными на промежутке $I$ коэффициентами и правой частью.
Для каждого $s\in I$ \textit{ядро Коши} $K(x,s)$ есть решение ЛОУ
$L_x[K]=0$ по переменной $x$ с условиями на диагонали
\[
 \partial_x^r K(s,s)=0\quad(0\le r\le n-2),\qquad
 \partial_x^{n-1}K(s,s)=1 .
\]
При $n=1$ остаётся только условие $K(s,s)=1$.
\end{definition}
\begin{result}{Интегральная формула Коши}
Для любого $\gamma\in I$ частное решение с нулевыми начальными данными
$Y^{(r)}(\gamma)=0$, $0\le r<n$, задаёт \textit{интегральная формула Коши}
\[
 Y(x)=\int_{\gamma}^{x}K(x,s)f(s)\,ds .
\]
\end{result}

\idea{Условия на диагонали обнуляют граничные члены до порядка $n-1$, а в $n$-й производной дают $f(x)$. Интегральная часть $L[Y]$ исчезает, поскольку $L_x[K]=0$.}

\end{document}
